\documentclass[11pt]{article}
\usepackage[a4paper,margin=25mm]{geometry}
\usepackage{amsmath,amssymb,booktabs,array,hyperref}
\hypersetup{colorlinks=true,urlcolor=blue,linkcolor=blue,pdftitle={Aggregate information and local-to-system sensitivity}}
\setlength{\parindent}{0pt}
\setlength{\parskip}{5pt}
\newcommand{\R}{\mathbb R}
\newcommand{\dd}{\mathrm d}
\title{Aggregate information and local-to-system sensitivity}
\author{Urban methodology review: mathematical companion to feedback 5}
\date{26 October 2026}
\begin{document}
\maketitle

\textbf{Traffic meaning.} First ask how a service reallocation changes discharge and the next traffic state. Then carry those changes through later access, queues and source waiting into complete total system time (TST). No packet-by-packet path differentiation is required. An aggregate rival may read the original fine-grained series through its declared macro conversion; report its retained representation and observation schedule.

This note clarifies the agreed principles, including the subsequent permission for original-series access through a model-defined macro boundary. Its algebra follows the accepted urban task's state-sensitivity contract. It does not select a missing access/service law, prove a useful gradient or construct a matched traffic pair.

\section*{1. From original records to macro state}

Let $D_{0:k}$ be the original fine-grained record available through time $k$. A rival $M$ may read that series and define its own aggregation or state estimator:
\begin{equation}
 r_k^{(M)}=\mathcal A_M(D_{0:k};\mathcal R_M).
\end{equation}
Here $\mathcal R_M$ describes its macro variables and resolution. Lane/bay masses, speeds, queue measures or estimated aggregate obstruction may be derived from the original records. The source can contain fine order; access to it does not itself make the model non-macro. We do not require a common projection to remove that information before every model applies its own method.

Declare the conversion and the complete retained representation, including estimator memory. A transient calculation of a useful macro summary is allowed. If an exact ordered packet ledger is also retained for later queries, that mesoscopic state belongs to the method's representation account. Successful aggregate reconstruction of a relevant effect must not be excluded as leakage.

For a causal observation-updated model, a complete retained macro record can be written
\begin{equation}
 I_k^{(M)}=\bigl(r_{0:k}^{(M)},\mu_k^{(M)},u_{0:k-1},\xi_{\rm common}\bigr),
\end{equation}
where $\mu_k^{(M)}$ is retained estimator memory, $u$ is applied service and $\xi_{\rm common}$ contains common physical/exogenous data. A model may use a subset. Full-record retrospective analysis is also legitimate; state when observations are available rather than present it as a forecast made before those outcomes. Macro state updated from the series is an allowed setup. Prediction without later updates is a particular capability, not a universal eligibility rule.

\textbf{A conditional information limit.} Fix one declared representation and observation protocol. Suppose two admissible full situations $a,b$ have the same complete retained record, $I^a=I^b=I$, for the response under examination. At the same action $u$, a predictor $h$ using only that record must give
\begin{equation}
 \widehat z^a=h(I,u)=\widehat z^b.
\end{equation}
Let their reference scalar responses be $z^a,z^b$. The triangle inequality gives
\begin{equation}
 \max\bigl\{|h(I,u)-z^a|,\ |h(I,u)-z^b|\bigr\}
 \ \geq\ \tfrac12 |z^a-z^b|.
 \label{eq:limit}
\end{equation}
Thus a response contrast exceeding twice a declared absolute error margin prevents both predictions from meeting that margin, independently of the chosen deterministic aggregate equation. The scalar response may be current discharge, a selected next-state measure, or a properly defined directional control response.

For a derivative comparison, hold the common predecision record fixed while changing the control. A differentiable $h(I,u)$ then supplies the same directional derivative to both situations. Equation~\eqref{eq:limit} applies to their reference directional responses as scalars. Equality of observations only at one nominal action does not justify equating derivatives of an observation-assisted, control-dependent record; those additional observation sensitivities must not be silently supplied.

This bound is conditional on the declared record. Equality under one aggregation does not establish equality under another model's legitimate summary of the original series. A universal result would require the condition across the admitted representation class; it has not been proved here. No matching pair has been constructed. Successful history-based reconstruction and failure of a selected subset-using model must be interpreted accordingly.

\textbf{Prediction and sensitivity meaning.} Local replay, observation-updated prediction and evolution without later updates are all legitimate tests. Disclose the conversion and observation schedule. Holding an observed record fixed gives a conditional response. If data or their estimated macro state change with the candidate control, their responses must be accounted for in a claimed total derivative. Supplying later outcomes does not establish that they were predicted before observation, but does not make the receiving representation non-macro. Matching a forward value at one action still does not establish matching sensitivity or richer-traffic validity.

\section*{2. One step: discharge response and next-state response}

Fix the starting state, exogenous data and a locally valid structural/active signature. For a \emph{selected} law write
\begin{equation}
 x_{k+1}=F_k(x_k,u_k),\qquad y_k=Y_k(x_k,u_k),
 \label{eq:transition}
\end{equation}
where $y_k$ is discharged continuous mass during the interval. Service fractions are dimensionless; $y_k$ is vehicle-equivalent mass. The continuous state includes the quantities and history endpoints that actually depend on the control. Structural labels are not differentiation variables.

Define the local partial responses
\begin{equation}
 A_k=F_{k,x},\quad B_k=F_{k,u},\qquad
 C_k=Y_{k,x},\quad D_k=Y_{k,u}.
\end{equation}
The first-order changes are
\begin{equation}
 \delta x_{k+1}=A_k\delta x_k+B_k\delta u_k,
 \qquad
 \delta y_k=C_k\delta x_k+D_k\delta u_k.
 \label{eq:local}
\end{equation}
Changing service \emph{now} at a fixed starting state means $\delta x_k=0$. The quantities of interest are then $D_k\delta u_k$ and $B_k\delta u_k$. They answer how current discharge and the following state change. If $F_k$ is defined by a joint allocation/equilibrium, its actual local response must be derived; merely writing~\eqref{eq:local} does not supply it.

For an illustrative \emph{movement-exclusive} service law from the neutral task, hold the available demand $E$ and receiving $R$ fixed:
\begin{equation}
 y=\min\{E,R,s a^\top u\,\Delta t\}.
\end{equation}
On a strict service-limited branch, a feasible direction $d$ gives
\begin{equation}
 \dd y[d]=s\Delta t\,a^\top d.
\end{equation}
On a strict demand/receiving-limited branch it gives zero. A queued movement can have a large service response when it has reachable demand and room. Conversely, blockage or jammed receiving can make additional service ineffective. This neutral illustration does not choose the unresolved shared-lane/bay law. At ties/front changes, a classical two-sided derivative need not exist.

\section*{3. A change last step, and a fixed allocation over time}

With $x_{k-1}$ held fixed, an isolated change to the previous control gives
\begin{equation}
 \delta x_k=B_{k-1}\delta u_{k-1},\qquad
 \delta y_k=C_kB_{k-1}\delta u_{k-1}
 \quad(\delta u_k=0).
 \label{eq:laststep}
\end{equation}
This is the precise form of ``what would be different this step if I had changed the knob last step?'' A persistent change adds the current direct term in~\eqref{eq:local} as well. Repeated application carries the response through later access, residual order, storage and source waiting.

In the accepted task, $g\in\R^4$ is one fixed physical allocation. A feasible reallocation direction $d$ satisfies $\mathbf1^\top d=0$ and respects the local lower bounds. It changes service share, not absolute cycle length. Write the applied allocation schedule as
\begin{equation}
 u_k(g)=\begin{cases}g,&k\text{ in the control window},\\
 \bar u_k,&k\text{ in the fixed clearance continuation}.
 \end{cases}
 \label{eq:control}
\end{equation}
Equation~\eqref{eq:control} makes explicit the retained reading of the existing D2 default: clearance service is independent of the candidate allocation. Its numerical schedule is not selected here. Initial preparation is common and independent of $g$.

For $S_k=\partial x_k/\partial g$ and $H_k=\partial u_k/\partial g$,
\begin{equation}
 S_0=0,\qquad S_{k+1}=A_kS_k+B_kH_k,
 \qquad H_k=\begin{cases}I_4,&\text{control},\\0,&\text{fixed clearance}.
 \end{cases}
 \label{eq:sequence}
\end{equation}
The state response generally persists during clearance through $A_kS_k$, even though the direct forcing $B_kH_k$ has ended. This is differentiation of a whole-system \emph{state sequence}, not a requirement to derive and sum individual packet paths.

\section*{4. Why the direct current TST charge can be unchanged}

Use a decision-boundary state after current exogenous arrivals have joined source availability. Under the inherited left-interval convention,
\begin{equation}
 \ell_k(x_k)=\Delta t\,[N_k^{\rm src}+N_k^{\rm net}],\qquad
 J(g)=\sum_{k=0}^{K-1}\ell_k(x_k(g)).
 \label{eq:tst}
\end{equation}
Network mass includes moving and queued traffic, counted once. The charge $\ell_k$ is fixed when $x_k$ is fixed, so its direct current-control partial is zero. Nevertheless, $B_k$ and $D_k$ may be nonzero, altering later state/time charges. Internal discharge between links also need not complete a trip. A valid zero local charge response therefore does not imply a zero complete-objective response.

For a common fixed horizon covering complete clearance in the relevant local neighborhood,
\begin{equation}
 \nabla_gJ=\sum_{k=0}^{K-1}S_k^\top\nabla_x\ell_k.
 \label{eq:tstgradient}
\end{equation}
Equivalently, define a future-time weight $\lambda_k$:
\begin{equation}
 \lambda_K=0,\qquad
 \lambda_k=\nabla_x\ell_k+A_k^\top\lambda_{k+1},\qquad
 \nabla_{u_k}J=B_k^\top\lambda_{k+1}.
 \label{eq:adjoint}
\end{equation}
The last expression values one local control change by its later time consequences. Summing $H_k^\top\nabla_{u_k}J$ gives the response to the shared fixed allocation. Summing only same-step discharge responses would miss those consequences.

Feedback 5 agrees the local-to-system principle. It does not select a horizon or endpoint algorithm. The formulas above use a common horizon under the stated clearance condition. A sound calculation through actual completion can evaluate the same objective. Nonclearing runs remain failed/censored. Do not add a positive mass cutoff or detach a state at clearance. A changing constraint or traffic front need not be smooth. Numerical correctness and useful richer-evaluator directions still require separate evidence.

\section*{Source and interpretation limits}

The state-transition, physical-control and complete-TST definitions follow the accepted urban task in \path{synthesis/experiment_formulation_formal.tex}, especially its derivative contract and full gradient path. The local/previous-step distinction, observation-bound argument and future-time interpretation above are algebraic explanations of those definitions, not newly selected traffic equations. The objective convention is also recorded in the approved M0 objective section and \path{src/diff_minimal_meso/objectives.py}.

Aggregate blockage capability is already represented in the reviewed literature: Liu and Chang's lane-group model addresses queue/blockage interactions; Wright et al. supply a macroscopic partial-FIFO restriction mechanism~\cite{liu,wright}. The bounded primary-source check supports that capability correction only. It does not establish exact order-conditioned gradient agreement, nor adopt either complete model here.

\begin{thebibliography}{2}
\bibitem{liu} Y. Liu and G.-L. Chang, ``An arterial signal optimization model for intersections experiencing queue spillback and lane blockage,'' \emph{Transp. Res. C}, vol. 19, pp. 130--144, 2011. doi: \href{https://doi.org/10.1016/j.trc.2010.04.005}{10.1016/j.trc.2010.04.005}. Primary check: abstract/introduction in \href{https://attap.umd.edu/wp-content/uploads/2022/05/TRC522.pdf}{author-hosted copy}.
\bibitem{wright} M. A. Wright, G. Gomes, R. Horowitz, and A. A. Kurzhanskiy, ``On node models for high-dimensional road networks,'' \emph{Transp. Res. B}, vol. 105, pp. 212--234, 2017. doi: \href{https://doi.org/10.1016/j.trb.2017.09.001}{10.1016/j.trb.2017.09.001}. Primary check: \href{https://arxiv.org/abs/1601.01054v3}{arXiv:1601.01054v3}, abstract.
\end{thebibliography}
\end{document}
